\textbf{Phần 1. Trắc nghiệm lựa chọn}
\setcounterpageref{ex}{0}

\begin{ex} %[2D14V73]
Một ngày có bao nhiêu giờ?
\choice
{$12$ giờ}
{\True $24$ giờ}
{$20$ giờ}
{$60$ giờ}
\loigiai{
Một ngày có đúng 24 giờ.
\textbf{Đáp án của câu hỏi là B}
}
\end{ex}

\begin{ex} %[2D14V73]
Tháng nào sau đây có đúng $30$ ngày?
\choice
{\True Tháng 4}
{Tháng 1}
{Tháng 3}
{Tháng 5}
\loigiai{
Các tháng có 30 ngày: Tháng 4, 6, 9, 11.
\textbf{Đáp án của câu hỏi là A}
}
\end{ex}

\begin{ex} %[2D14V73]
Đổi: $1\text{ m } 8\text{ dm} = \dots\text{ dm}$.
\choice
{$108\text{ dm}$}
{$81\text{ dm}$}
{\True $18\text{ dm}$}
{$9\text{ dm}$}
\loigiai{
$1\text{ m} = 10\text{ dm} \Rightarrow 1\text{ m } 8\text{ dm} = 18$ dm.
\textbf{Đáp án của câu hỏi là C}
}
\end{ex}

\begin{ex} %[2D14V73]
Mẹ mua một gói mì chính cân nặng $500\text{ g}$ và một gói đường cân nặng $500\text{ g}$. Cả hai gói cân nặng bao nhiêu ki-lô-gam?
\choice
{$2\text{ kg}$}
{$10\text{ kg}$}
{$5\text{ kg}$}
{\True $1\text{ kg}$}
\loigiai{
$500\text{ g} + 500\text{ g} = 1 000\text{ g} = 1\text{ kg}$.
\textbf{Đáp án của câu hỏi là D}
}
\end{ex}

\begin{ex} %[2D14V73]
Đồng hồ chỉ 8 giờ 15 phút. Kim phút đang chỉ vào số nào?
\choice
{Số 8}
{\True Số 3}
{Số 6}
{Số 12}
\loigiai{
Kim phút chỉ số 3 là 15 phút.
\textbf{Đáp án của câu hỏi là B}
}
\end{ex}

\begin{ex} %[2D14V73]
Có 3 tờ 200 đồng và 1 tờ 500 đồng. Tổng số tiền là
\choice
{\True $1 100\text{ đồng}$}
{$700\text{ đồng}$}
{$900\text{ đồng}$}
{$1 000\text{ đồng}$}
\loigiai{
$200 \times 3 + 500 = 600 + 500 = 1 100$ đồng.
\textbf{Đáp án của câu hỏi là A}
}
\end{ex}

\begin{ex} %[2D14V73]
Một can chứa 10 lít dầu. Người ta rót vào 2 can nhỏ, mỗi can 4 lít. Trong can còn lại bao nhiêu lít dầu?
\choice
{$3$ lít}
{$4$ lít}
{\True $2$ lít}
{$1$ lít}
\loigiai{
Đã rót: $4 \times 2 = 8$ l. Còn lại: $10 - 8 = 2$ lít.
\textbf{Đáp án của câu hỏi là C}
}
\end{ex}

\begin{ex} %[2D14V73]
Điền dấu thích hợp: $1\text{ km} \dots 950\text{ m} + 50\text{ m}$.
\choice
{$>$}
{$<$}
{Không so sánh được}
{\True $=$}
\loigiai{
$950\text{ m} + 50\text{ m} = 1 000\text{ m} = 1\text{ km}$. Hai vế bằng nhau.
\textbf{Đáp án của câu hỏi là D}
}
\end{ex}

\begin{ex} %[2D14V73]
Hôm nay là ngày 20 tháng 5. Còn bao nhiêu ngày nữa thì hết tháng 5?
\choice
{$10$ ngày}
{\True $11$ ngày}
{$12$ ngày}
{$9$ ngày}
\loigiai{
Tháng 5 có 31 ngày. Còn: $31 - 20 = 11$ ngày.
\textbf{Đáp án của câu hỏi là B}
}
\end{ex}

\begin{ex} %[2D14V73]
Nam bắt đầu học bài lúc 8 giờ tối và học xong lúc 8 giờ 45 phút tối. Nam đã học trong thời gian bao lâu?
\choice
{\True $45$ phút}
{$40$ phút}
{$50$ phút}
{$1$ giờ}
\loigiai{
Thời gian học là 45 phút.
\textbf{Đáp án của câu hỏi là A}
}
\end{ex}

\begin{ex} %[2D14V73]
Con gà nặng $2\text{ kg}$, con ngỗng nặng hơn con gà $3\text{ kg}$. Cả hai con cân nặng bao nhiêu ki-lô-gam?
\choice
{$5\text{ kg}$}
{$6\text{ kg}$}
{\True $7\text{ kg}$}
{$8\text{ kg}$}
\loigiai{
Con ngỗng: $2 + 3 = 5$ kg. Cả hai con: $2 + 5 = 7$ kg.
\textbf{Đáp án của câu hỏi là C}
}
\end{ex}

\begin{ex} %[2D14V73]
Để có đúng $1 000\text{ đồng}$ từ các tờ tiền $200\text{ đồng}$, ta cần bao nhiêu tờ tiền?
\choice
{$4$ tờ}
{$6$ tờ}
{$10$ tờ}
{\True $5$ tờ}
\loigiai{
$1 000 : 200 = 5$ tờ.
\textbf{Đáp án của câu hỏi là D}
}
\end{ex}

\textbf{Phần 2. Trắc nghiệm đúng sai}
\setcounterpageref{ex}{0}

\begin{ex} %[2D14V73]
Xét tính đúng sai của các đơn vị và số đo đại lượng:
\choiceTF
{\True 1 m = 10 dm = 100 cm.}
{Tháng 2 luôn luôn có 30 ngày.}
{\True 1 kg nặng hơn 900 g.}
{\True 5 tờ 100 đồng có giá trị bằng 1 tờ 500 đồng.}
\loigiai{
\begin{itemize}
    \item Ý a) Đúng: Đúng.
    \item Ý b) Sai: Tháng 2 chỉ có 28 hoặc 29 ngày. Sai.
    \item Ý c) Đúng: 1 kg = 1 000 g > 900 g. Đúng.
    \item Ý d) Đúng: Đúng, đều bằng 500 đồng.
\end{itemize}
\begin{itemize}
    \item a) \textbf{Đúng}.
    \item b) \textbf{Sai}.
    \item c) \textbf{Đúng}.
    \item d) \textbf{Đúng}.
\end{itemize}
}
\end{ex}

\begin{ex} %[2D14V73]
Một chuyến xe buýt khởi hành lúc 7 giờ sáng và đến nơi lúc 8 giờ 30 phút sáng cùng ngày:
\choiceTF
{\True Chuyến đi kéo dài 1 giờ 30 phút.}
{\True 1 giờ 30 phút bằng 90 phút.}
{Xe đến nơi vào buổi chiều.}
{\True Nếu xe khởi hành muộn 15 phút thì đến nơi lúc 8 giờ 45 phút.}
\loigiai{
\begin{itemize}
    \item Ý a) Đúng: Đúng.
    \item Ý b) Đúng: 60 + 30 = 90 phút. Đúng.
    \item Ý c) Sai: Xe đến lúc 8 giờ 30 sáng. Sai.
    \item Ý d) Đúng: Đúng.
\end{itemize}
\begin{itemize}
    \item a) \textbf{Đúng}.
    \item b) \textbf{Đúng}.
    \item c) \textbf{Sai}.
    \item d) \textbf{Đúng}.
\end{itemize}
}
\end{ex}

\textbf{Phần 3. Trả lời ngắn}
\setcounterpageref{ex}{0}

\begin{ex} %[2D14V73]
Một thùng sơn nặng 18 kg. Bác thợ dùng hết 5 kg. Thùng sơn còn lại nặng bao nhiêu ki-lô-gam?
\shortans{13}
\loigiai{
\textbf{Đáp số:} $13$\\
$18 - 5 = 13$ kg.
}
\end{ex}

\begin{ex} %[2D14V73]
Một bể nước chứa 50 lít nước. Sau khi dùng hết 32 lít, người ta đổ thêm vào 15 lít nữa. Trong bể có bao nhiêu lít nước?
\shortans{33}
\loigiai{
\textbf{Đáp số:} $33$\\
$50 - 32 + 15 = 33$ lít.
}
\end{ex}

\begin{ex} %[2D14V73]
Bác Nam có 1 tờ 1 000 đồng, bác mua vé xe buýt hết 700 đồng. Bác còn lại bao nhiêu đồng?
\shortans{300}
\loigiai{
\textbf{Đáp số:} $300$\\
$1 000 - 700 = 300$ đồng.
}
\end{ex}

\begin{ex} %[2D14V73]
Đoạn đường từ nhà đến trường dài 500 m, Nam đi học cả đi và về hết bao nhiêu ki-lô-mét?
\shortans{1}
\loigiai{
\textbf{Đáp số:} $1$\\
$500 + 500 = 1 000\text{ m} = 1$ km.
}
\end{ex}

\begin{ex} %[2D14V73]
Tháng 12 có bao nhiêu ngày?
\shortans{31}
\loigiai{
\textbf{Đáp số:} $31$\\
Tháng 12 có 31 ngày.
}
\end{ex}

\begin{ex} %[2D14V73]
Một hình vuông có cạnh dài 2 dm. Chu vi hình vuông đó là bao nhiêu xăng-ti-mét?
\shortans{80}
\loigiai{
\textbf{Đáp số:} $80$\\
Cạnh = 20 cm. Chu vi: $20 \times 4 = 80$ cm.
}
\end{ex}
